\chapter{Process Mining Canon as Type Law}
\label{ch:processcanon}

The process mining literature has, over four decades, accumulated a rich set of formal
objects: event logs, Petri nets, workflow nets, process trees, conformance metrics, object
lifecycles, and their interrelations. These objects are defined by mathematical structures
in papers --- not as runtime values, but as sets, relations, partial orders, and typed
functions with named laws.

This chapter demonstrates how each canonical formal object from that literature is encoded
in \texttt{wasm4pm-compat} as a zero-cost Rust type. The encoding is not approximate: the
compile-time constraints of the Rust type system enforce the structural laws from the
papers. Out-of-law configurations are rejected by the compiler, not caught by a test or a
runtime guard.

% -----------------------------------------------------------------------------
\section{OCEL 2.0 Formal Objects}
\label{sec:canon-ocel}
% -----------------------------------------------------------------------------

The Object-Centric Event Log 2.0 standard~\cite{ghahfarokhi2021ocel,berti2023ocel20}
replaces the classical single-case-notion event log with a multi-relational graph whose
nodes are typed objects and timestamped events, and whose edges are qualified links. This
is not a minor extension: OCEL 2.0 is a fundamentally different metamodel that breaks the
classical assumption that every event belongs to exactly one case.

\subsection{The OCEL 2.0 Metamodel}

The formal metamodel of OCEL 2.0 is a bipartite, typed graph over events and objects,
extended with object-to-object relations, object attribute evolution, and qualified link
roles. Figure~\ref{fig:ocel-metamodel} shows the metamodel as a TikZ diagram.

\begin{figure}[ht]
\centering
\begin{tikzpicture}[
  node distance=3.5cm and 4.5cm,
  every node/.style={font=\small},
  boxnode/.style={draw, rectangle, rounded corners=2pt, minimum height=1.1cm,
                  minimum width=2.5cm, align=center, fill=white},
  linknode/.style={draw, rectangle, rounded corners=2pt, minimum height=0.9cm,
                   minimum width=3.2cm, align=center, fill=gray!10, font=\small\itshape},
  changesnode/.style={draw, rectangle, rounded corners=2pt, minimum height=0.9cm,
                       minimum width=2.8cm, align=center, fill=gray!10, font=\small\itshape},
  >=latex
]

% Main nodes
\node[boxnode] (event)  {$\texttt{OcelEvent}$\\{\tiny id, activity, timestamp, attrs}};
\node[boxnode, right=of event] (object) {$\texttt{OcelObject}$\\{\tiny id, type, attrs}};

% Link types
\node[linknode, below=1.5cm of event, xshift=1.8cm] (e2o) {$\texttt{EventObjectLink}$\\(E2O)\\{\tiny qualifier: Option<String>}};

\node[linknode, below=1.5cm of object, xshift=0.6cm] (o2o) {$\texttt{ObjectObjectLink}$\\(O2O)\\{\tiny qualifier: Option<String>}};

% Object change
\node[changesnode, below=3.2cm of object] (change)
  {$\texttt{ObjectChange}$\\{\tiny object\_id, attribute, value, ts}};

% Container
\node[boxnode, above=1.8cm of event, xshift=2cm] (log) {$\texttt{OcelLog}$\\{\tiny events, objects, e2o, o2o, changes}};

% Edges: log contains events and objects
\draw[->] (log) -- node[above left, font=\scriptsize] {events} (event);
\draw[->] (log) -- node[above right, font=\scriptsize] {objects} (object);

% Edges: E2O link
\draw[->, dashed] (event) -- node[right, font=\scriptsize] {event\_id} (e2o);
\draw[->, dashed] (object) -- node[left, font=\scriptsize] {object\_id} (e2o);

% Edges: O2O link
\draw[->, dashed] (object) to[bend left=15] node[above, font=\scriptsize] {source\_id} (o2o);
\draw[->, dashed] (object) to[bend right=15] node[below, font=\scriptsize] {target\_id} (o2o);

% Edges: ObjectChange
\draw[->, dashed] (object) -- node[right, font=\scriptsize] {object\_id} (change);

\end{tikzpicture}
\caption{The OCEL 2.0 metamodel as encoded in \texttt{wasm4pm-compat}. Every node is a
  zero-cost Rust type. \texttt{EventObjectLink} carries a qualified E2O edge; dangling
  links are refused at \texttt{OcelLog::validate} time.}
\label{fig:ocel-metamodel}
\end{figure}

\subsection{Zero-Cost Encoding of the Metamodel}

The OCEL 2.0 metamodel is encoded in \texttt{src/ocel.rs} using five mutually distinct
Rust types.

\begin{definition}[OCEL 2.0 Formal Objects]
\label{def:ocel-formal}
Let $\mathcal{L}$ be an OCEL 2.0 log. The OCEL metamodel defines:
\begin{itemize}
  \item $\mathcal{E}$ --- a set of events, each carrying an activity name, optional
        timestamp, and typed attributes.
  \item $\mathcal{O}$ --- a set of typed objects, each with a stable identifier,
        an object type, and typed attributes.
  \item $\text{E2O} \subseteq \mathcal{E} \times \mathcal{O} \times \mathcal{Q}$ ---
        event-to-object links, where $\mathcal{Q}$ is a (possibly empty) qualifier set.
  \item $\text{O2O} \subseteq \mathcal{O} \times \mathcal{O} \times \mathcal{Q}$ ---
        object-to-object links with optional qualifiers.
  \item $\Delta \subseteq \mathcal{O} \times \text{Attr} \times \text{Val} \times
        \mathbb{T}$ --- timestamped attribute-evolution records.
\end{itemize}
A log $\mathcal{L} = (\mathcal{E}, \mathcal{O}, \text{E2O}, \text{O2O}, \Delta)$ is
\emph{structurally valid} iff every link in E2O and O2O references a declared entity.
\end{definition}

\begin{lstlisting}[language=Rust,
  caption={The five OCEL 2.0 formal objects as Rust types. Structure-only --- no engine
    logic, no discovery.},
  label={lst:ocel-types}]
// An OCEL event: id, activity, optional timestamp, typed attributes.
pub struct OcelEvent {
    id:           String,
    activity:     String,
    timestamp_ns: Option<i64>,
    attributes:   Vec<OcelAttribute>,
}

// An OCEL object: stable id, typed identity, initial attribute set.
pub struct OcelObject {
    id:          String,
    object_type: String,
    attributes:  Vec<OcelAttribute>,
}

// E2O link: typed edge with optional role qualifier.
pub struct EventObjectLink {
    event_id:  String,
    object_id: String,
    qualifier: Option<String>,
}

// O2O link: typed relationship between two objects.
pub struct ObjectObjectLink {
    source_id: String,
    target_id: String,
    qualifier: Option<String>,
}

// Attribute evolution: object-id.attr = value at timestamp.
pub struct ObjectChange {
    object_id:    String,
    attribute:    String,
    value:        String,
    timestamp_ns: Option<i64>,
}

// Container: the OCEL 2.0 log as a first-class type.
pub struct OcelLog {
    events:  Vec<OcelEvent>,
    objects: Vec<OcelObject>,
    e2o:     Vec<EventObjectLink>,
    o2o:     Vec<ObjectObjectLink>,
    changes: Vec<ObjectChange>,
}
\end{lstlisting}

\subsection{OcelDims: Dimensional Schema Law}

A real OCEL log does not carry a random bag of objects: it has a declared schema of object
types and activity types, with constraints on which activities may touch which object types.
The \texttt{OcelDims} type in \texttt{src/ocel.rs} encodes this dimensional schema.

\begin{definition}[OCEL Dimensional Schema]
\label{def:ocel-dims}
An OCEL dimensional schema is a triple $(\mathcal{OT}, \mathcal{AT}, \lambda)$ where
$\mathcal{OT}$ is a set of object types, $\mathcal{AT}$ is a set of activity types, and
$\lambda: \mathcal{AT} \to 2^{\mathcal{OT}}$ maps each activity to the set of object types
it is permitted to reference. An E2O link $(e, o, q)$ is \emph{type-lawful} iff
$\text{type}(o) \in \lambda(\text{activity}(e))$.
\end{definition}

The structural integrity of \texttt{OcelLog} is enforced by \texttt{OcelLog::validate},
which refuses dangling links as \texttt{OcelRefusal::DanglingEventObjectLink} and
\texttt{OcelRefusal::DanglingObjectObjectLink}. These refusals carry named law reasons --- not
a bare \texttt{InvalidInput} string, but a specific variant that names the violated
structural invariant. The audit script \texttt{audit\_no\_unsealed\_refusal.sh} enforces
this: no refusal variant may use a catch-all string.

% -----------------------------------------------------------------------------
\section{XES and Case-Centric Evidence}
\label{sec:canon-xes}
% -----------------------------------------------------------------------------

The XES standard (IEEE 1849)~\cite{verbeek2010xes,ieee2016xes} defines a case-centric
event log interchange format: a log contains traces, traces contain events, and every event
belongs to exactly one trace (case). XES is fundamentally different from OCEL 2.0: the
single case notion is not a limitation to be worked around but a structural invariant that
the format enforces.

\begin{definition}[XES Type Separation]
\label{def:xes-separation}
\texttt{XesLog} and \texttt{OcelLog} are distinct types in \texttt{wasm4pm-compat}. There
is no implicit coercion, no common base type, and no shared constructor. A function
expecting a \texttt{XesLog} cannot receive an \texttt{OcelLog}; the compiler rejects the
substitution as a type mismatch.
\end{definition}

\subsection{The CaseCentricMarker Boundary Type}

The XES module introduces \texttt{CaseCentricMarker}: a zero-sized type that marks any
type or value as carrying case-centric (not object-centric) semantics.

\begin{lstlisting}[language=Rust,
  caption={\texttt{CaseCentricMarker}: a zero-sized boundary type making the case-centric
    law visible in the type system. Carrying this marker in a \texttt{PhantomData} field
    makes case-centric evidence incompatible with object-centric slots.},
  label={lst:case-centric-marker}]
/// A zero-sized marker that tags a type or value as case-centric.
///
/// Use as PhantomData<CaseCentricMarker> to prevent OcelLog substitution
/// for XesLog at the type level.
#[derive(Clone, Copy, Debug, PartialEq, Eq, Hash, Default)]
pub struct CaseCentricMarker;
\end{lstlisting}

\subsection{XES Extension Prefix Law}

XES attributes are namespaced by declared extensions (e.g., \texttt{concept:name},
\texttt{time:timestamp}). An \texttt{XesExtension} with an empty prefix is structurally
invalid: the prefix is the namespace separator, and an empty prefix makes all attributes
globally unnamespaced. The \texttt{XesRefusal::InvalidExtension} variant names this
violation.

\begin{lstlisting}[language=Rust,
  caption={XES extension prefix law: an empty prefix is refused. The law is named
    \texttt{XesRefusal::InvalidExtension}, not a bare string error.},
  label={lst:xes-extension}]
pub struct XesExtension {
    name:   String,
    prefix: String,   // must be non-empty; refuses as InvalidExtension
    uri:    String,
}

pub enum XesRefusal {
    InvalidExtension,  // prefix is empty
    MissingConceptName,
    MissingTimestamp,
}
\end{lstlisting}

\subsection{Compile-Fail Receipt: XesLog Cannot Substitute for OcelLog}

The ALIVE gate for the XES/OCEL boundary is enforced by two compile-fail receipts. The
first demonstrates that \texttt{OcelLog} cannot substitute for \texttt{XesLog}
(Listing~\ref{lst:ocel-as-xes-fail}); the second demonstrates the reverse direction.

\begin{lstlisting}[language=Rust,
  caption={Compile-fail receipt: \texttt{ocel\_log\_as\_xes\_log.rs}.
    Law: OcelLog and XesLog are distinct types;
    object-centric logs cannot replace case-centric logs.},
  label={lst:ocel-as-xes-fail}]
// COMPILE-FAIL: OCEL/XES confusion law
// Law: OcelLog and XesLog are structurally distinct types.
// Object-centric logs cannot silently replace case-centric logs.
use wasm4pm_compat::ocel::OcelLog;
use wasm4pm_compat::xes::XesLog;

fn requires_xes_log(_log: XesLog) {}

fn main() {
    let ocel: OcelLog = todo!();
    requires_xes_log(ocel);  // ERROR: expected XesLog, found OcelLog
}
\end{lstlisting}

The compiler error for this fixture is:

\begin{verbatim}
error[E0308]: mismatched types
  --> tests/ui/compile_fail/ocel_log_as_xes_log.rs:12:22
   |
12 |     requires_xes_log(ocel);
   |                      ^^^^ expected `XesLog`, found `OcelLog`
\end{verbatim}

This error is the type-law receipt for the OCEL/XES boundary law. A
fixture that fails for a different reason --- a missing import, a typo in the function name
--- is not a valid receipt for this law. The exact \texttt{.stderr} snapshot paired with the
fixture proves that the failure is for the intended law.

% -----------------------------------------------------------------------------
\section{Petri Net Bipartite Arc Law}
\label{sec:canon-petri}
% -----------------------------------------------------------------------------

The Petri net bipartite arc law is one of the oldest structural invariants in process
mining. Murata~\cite{murata1989petri} defines a Petri net as a bipartite directed graph
over places and transitions, where arcs are allowed only from places to transitions (the
pre-incidence relation $W^-$) or from transitions to places (the post-incidence relation
$W^+$). No arc from a place to a place, or from a transition to a transition, may exist.

\begin{definition}[Petri Net Bipartite Arc Law (Murata 1989)]
\label{def:petri-bipartite}
A Petri net is a tuple $N = (P, T, F, W, M_0)$ where $P$ and $T$ are disjoint finite
sets of places and transitions, $F \subseteq (P \times T) \cup (T \times P)$ is the flow
relation (bipartite arcs only), $W: F \to \mathbb{N}^+$ is the arc weight function, and
$M_0: P \to \mathbb{N}$ is the initial marking. The constraint
$F \cap (P \times P) = \emptyset$ and $F \cap (T \times T) = \emptyset$ is the
\emph{bipartite arc law}.
\end{definition}

\subsection{Bipartite Law as a Sealed Trait}

The bipartite arc law is enforced by making place-to-place and transition-to-transition
arcs \emph{unconstructible}. The only arc types in \texttt{src/petri.rs} are
\texttt{PlaceToTransitionArc<P, T, Weight>} and \texttt{TransitionToPlaceArc<T, P, Weight>}.
There is no \texttt{PlaceToPlaceArc} type. The type system guarantees that no such arc
can ever be constructed.

\begin{lstlisting}[language=Rust,
  caption={\texttt{PlaceToTransitionArc} and \texttt{TransitionToPlaceArc}: the only two
    arc inhabitants. No \texttt{PlaceToPlaceArc} type exists in this crate.},
  label={lst:petri-arcs}]
/// Pre-incidence arc: place -> transition (W-).
/// Paper: Murata (1989) S.2 --- F <= (PxT) U (TxP).
pub struct PlaceToTransitionArc<P, T, Weight> {
    pub(crate) _from: PhantomData<P>,
    pub(crate) _to:   PhantomData<T>,
    pub weight: Weight,
}

/// Post-incidence arc: transition -> place (W+).
/// Paper: Murata (1989) S.2 --- bipartite arc law.
pub struct TransitionToPlaceArc<T, P, Weight> {
    pub(crate) _from: PhantomData<T>,
    pub(crate) _to:   PhantomData<P>,
    pub weight: Weight,
}
\end{lstlisting}

\begin{theorem}[Petri Net Arc Law Soundness]
\label{thm:petri-arc-law}
In a Petri net $N = (P, T, F)$, the flow relation $F \subseteq (P \times T) \cup (T \times P)$.
The type system of \texttt{wasm4pm-compat} enforces this by making
\texttt{PlaceToTransitionArc} and \texttt{TransitionToPlaceArc} the only inhabitants
of the bipartite arc universe: there is no type \texttt{PlaceToPlaceArc} or
\texttt{TransitionToTransitionArc}. All combinations excluded by Murata's definition are
rejected at compile time.
\end{theorem}

\begin{proof}
Let $\mathcal{T}$ be the set of types in \texttt{src/petri.rs} whose name contains
\texttt{Arc}. By inspection, $\mathcal{T} = \{\texttt{PlaceToTransitionArc},
\texttt{TransitionToPlaceArc}, \texttt{BipartiteArcConst}\}$. Each of these types either
requires distinct place and transition type parameters, or carries a const generic
\texttt{ArcDirectionConst} that distinguishes the two directions. The
\texttt{ArcDirectionConst} enum has exactly two variants:
\texttt{PlaceToTransition} and \texttt{TransitionToPlace}. A slot requiring
\texttt{PlaceToTransitionArc} cannot receive a \texttt{TransitionToPlaceArc}, as
demonstrated by the compile-fail receipt in Listing~\ref{lst:petri-place-place-fail}.
\end{proof}

\begin{lstlisting}[language=Rust,
  caption={Compile-fail receipt: \texttt{petri\_place\_to\_place\_arc.rs}.
    Law: BipartitePetriArcLaw --- P->P arcs are unconstructible (Murata 1989 S.2).},
  label={lst:petri-place-place-fail}]
// COMPILE-FAIL: Bipartite arc law --- place->place arcs are unconstructible.
// Paper: Murata (1989) S.2 --- F <= (PxT) U (TxP), no P->P arcs.
use wasm4pm_compat::petri::PlaceToTransitionArc;
use wasm4pm_compat::petri::TransitionToPlaceArc;

struct P1; struct P2;

fn main() {
    // Attempting to assign a TransitionToPlace arc into a PlaceToTransition slot:
    let arc: PlaceToTransitionArc<P1, P2, u8> =
        TransitionToPlaceArc::<P1, P2, u8>::new(1u8);
    // ERROR: expected PlaceToTransitionArc, found TransitionToPlaceArc
    let _ = arc;
}
\end{lstlisting}

The const-generic \texttt{BipartiteArcConst<DIR, Weight>} provides a single parameterized
arc type for containers that need to hold arcs of either direction while still encoding
direction in the type. \texttt{BipartiteArcConst<\{PlaceToTransition\}, u8>} and
\texttt{BipartiteArcConst<\{TransitionToPlace\}, u8>} are distinct types at compile time:
a slot requiring the pre-incidence direction rejects the post-incidence direction as a
type error.

% -----------------------------------------------------------------------------
\section{WF-Net Soundness}
\label{sec:canon-wfnet}
% -----------------------------------------------------------------------------

A workflow net (WF-net) is a Petri net with a single source place, a single sink place,
and the requirement that every node is on some path from source to sink. van der
Aalst~\cite{vanderaalst1998workflow} defines \emph{soundness} as the conjunction of three
properties: option to complete (every reachable marking can reach the final marking), proper
completion (reaching the final marking empties the net), and no dead transitions (every
transition is fireable from some reachable marking).

Computing soundness is decidable but non-trivial (it reduces to reachability in a
state space that may be exponential). This crate does not compute soundness: that is an
engine and graduates to \texttt{wasm4pm}. Instead, a WF-net carries a \emph{soundness
state} token at the type level, recording whether soundness is unknown, merely claimed by
an upstream assertion, or witnessed by the \texttt{wasm4pm} engine.

\begin{definition}[Non-Forgeable Soundness Witness]
\label{def:wfnet-soundness}
Let $\mathcal{S} = \{\texttt{Unknown}, \texttt{Claimed}, \texttt{Witnessed}\}$ be the
set of soundness states. A WF-net in \texttt{wasm4pm-compat} carries its soundness state
as a const generic parameter: \texttt{WfNetConst<SOUNDNESS>}. The \texttt{Witnessed}
variant is non-forgeable: the only public path to
\texttt{WfNetConst<\{SoundnessState::Witnessed\}>} is through the \texttt{witness\_soundness}
function, which requires a witness value of type produced by the \texttt{wasm4pm} engine.
\end{definition}

\begin{lstlisting}[language=Rust,
  caption={The three WF-net soundness states as a \texttt{ConstParamTy} enum in
    \texttt{src/law.rs}. \texttt{Witnessed} is the non-forgeable state.},
  label={lst:wfnet-soundness}]
#[derive(ConstParamTy, PartialEq, Eq, Clone, Copy, Debug, Hash)]
pub enum SoundnessState {
    /// Soundness not yet asserted (default state at construction).
    Unknown,
    /// Soundness claimed by upstream; unproven in this crate.
    Claimed,
    /// Soundness witnessed via wasm4pm engine output (non-forgeable).
    Witnessed,
}
\end{lstlisting}

The three soundness states model three distinct epistemic positions. A
\texttt{WfNetConst<\{Unknown\}>} is a WF-net whose soundness has not been assessed. A
\texttt{WfNetConst<\{Claimed\}>} carries an upstream assertion --- it may have been declared
sound by a modeling tool that did not verify it. A \texttt{WfNetConst<\{Witnessed\}>}
carries a proof from the \texttt{wasm4pm} engine that the three soundness properties hold.

\begin{remark}[Soundness is a claim, not a result]
\label{rem:soundness-claim}
This crate never promotes a net from \texttt{Unknown} or \texttt{Claimed} to
\texttt{Witnessed} by itself. Promoting requires calling \texttt{witness\_soundness}, which
is only accessible when the \texttt{wasm4pm} feature is enabled and which requires a
\texttt{SoundnessProof} produced by the engine. A function that requires a
\texttt{WfNetConst<\{Witnessed\}>} parameter cannot be called with
\texttt{WfNetConst<\{Unknown\}>}: the const generic parameter makes them distinct types.
\end{remark}

% -----------------------------------------------------------------------------
\section{POWL Projection Law}
\label{sec:canon-powl}
% -----------------------------------------------------------------------------

The Partially Ordered Workflow Language (POWL)~\cite{kourani2024powl} extends process
trees with partial orders, allowing models that can express concurrent behavior without
forcing a total order. A POWL model is a graph of nodes --- atoms, partial orders, exclusive
choices, loops, and silent steps --- with precedence edges that must form a DAG. The key
insight of POWL is that not every partial order can be represented as a block-structured
process tree: some POWL fragments are \emph{irreducible}, meaning that any process-tree
projection would be lossy.

\begin{definition}[Process-Tree Projectability (Kourani 2024)]
\label{def:powl-projectable}
A POWL model $M$ is \emph{process-tree projectable} iff every partial order node in $M$
can be expressed as a binary sequence, choice, or loop operator in a block-structured
process tree $T$ such that $\mathcal{L}(M) = \mathcal{L}(T)$ (language equivalence).
A POWL model containing at least one irreducible partial order \emph{exceeds the process
tree} --- projection would silently lose behavior.
\end{definition}

\subsection{The Sealed Assertion Pattern}

The projection law is enforced by two sealed marker types and the sealed
\texttt{TreeProjectable} trait. Only \texttt{ProcessTreeProjectable} satisfies
\texttt{TreeProjectable}; \texttt{ExceedsProcessTree} is deliberately excluded from the
seal. A function requiring a \texttt{TreeProjectable} argument cannot receive an
\texttt{ExceedsProcessTree} marker.

\begin{lstlisting}[language=Rust,
  caption={The POWL projection law as a sealed trait. Only \texttt{ProcessTreeProjectable}
    satisfies \texttt{TreeProjectable}. \texttt{ExceedsProcessTree} is excluded from the
    seal, making it a compile error to project a non-projectable model.},
  label={lst:powl-sealed}]
mod tree_projectable_seal {
    pub trait Sealed {}
    impl Sealed for super::ProcessTreeProjectable {}
    // ExceedsProcessTree deliberately NOT sealed here.
}

/// Sealed marker: only ProcessTreeProjectable satisfies this bound.
pub trait TreeProjectable: tree_projectable_seal::Sealed {}
impl TreeProjectable for ProcessTreeProjectable {}

/// Structural gate: only POWL markers that are tree-projectable pass.
/// This is NOT a discovery function --- it proves projection law at type level.
pub fn assert_tree_projectable<P: TreeProjectable>(_marker: P) -> bool {
    true
}
\end{lstlisting}

The \texttt{assert\_tree\_projectable} function is not a runtime assertion. It is a
type-level gate: calling it with \texttt{ExceedsProcessTree} is a compile error.

\begin{lstlisting}[language=Rust,
  caption={Compile-fail receipt for the POWL projection law.
    \texttt{ExceedsProcessTree} is not \texttt{TreeProjectable}.},
  label={lst:powl-exceeds-fail}]
// COMPILE-FAIL: POWL projection law
// ExceedsProcessTree is not TreeProjectable.
// Paper: Kourani et al. (2024) --- irreducible partial orders exceed process trees.
use wasm4pm_compat::powl::{assert_tree_projectable, ExceedsProcessTree};

fn main() {
    assert_tree_projectable(ExceedsProcessTree);
    // ERROR: ExceedsProcessTree does not implement TreeProjectable
}
\end{lstlisting}

\subsection{Acyclic Partial Order Witness}

A POWL partial order must be a DAG: the precedence relation must be irreflexive,
asymmetric, and transitive. The \texttt{AcyclicPartialOrder} witness marker records that
this assertion has been made. The sealed \texttt{AcyclicWitness} trait ensures that only
\texttt{AcyclicPartialOrder} --- not the bare \texttt{PartialOrder} shape marker --- satisfies
the acyclicity requirement.

\begin{remark}[Witness markers record assertions, not computations]
\label{rem:powl-witness}
Neither \texttt{assert\_tree\_projectable} nor \texttt{assert\_acyclic} runs an
algorithm. They record that an assertion was made at a specific site in the code. The
actual cycle-detection and process-tree equivalence check graduate to \texttt{wasm4pm}.
The type system then carries the assertion evidence into downstream consumers that require
it.
\end{remark}

% -----------------------------------------------------------------------------
\section{Conformance Metrics}
\label{sec:canon-conformance}
% -----------------------------------------------------------------------------

Carmona et al.~\cite{carmona2018conformance} define the four standard conformance quality
dimensions: \emph{fitness} (the model can reproduce the log behavior), \emph{precision}
(the model does not allow excessive behavior absent from the log), \emph{generalization}
(the model does not overfit), and \emph{simplicity} (the model is as simple as possible).
Each dimension is a real value in $[0, 1]$, where higher is better. An F1 score
aggregates fitness and precision. All five dimensions are bounded: a metric value outside
$[0, 1]$ is a structural defect.

\begin{definition}[Between01-Bounded Metric]
\label{def:between01}
A conformance metric value is a rational $\text{NUM} / \text{DEN} \in [0, 1]$. In
\texttt{wasm4pm-compat}, this is enforced by the
\texttt{Between01<\textrm{NUM, DEN}>} type: it requires $\text{DEN} > 0$ and
$\text{NUM} \leq \text{DEN}$ at the type level. A \texttt{Between01<2, 1>} does not
compile.
\end{definition}

\subsection{The Metric Type}

The \texttt{Metric<KIND, NUM, DEN>} type in \texttt{src/conformance.rs} carries a
conformance dimension (\texttt{QualityMetricKind}), a numerator, and a denominator. Both
the \texttt{DEN > 0} and \texttt{NUM <= DEN} constraints are enforced by
\texttt{Require\{...\}: IsTrue} where-bounds at the type level.

\begin{lstlisting}[language=Rust,
  caption={\texttt{Metric<KIND, NUM, DEN>}: conformance score with compile-time
    $[0,1]$ bounds. The \texttt{FitnessConst<2, 1>} compile-fail receipt proves that
    out-of-range scores are rejected.},
  label={lst:metric-type}]
pub struct Metric<const KIND: QualityMetricKind,
                  const NUM: u64,
                  const DEN: u64>
where
    Require<{ DEN > 0 }>:   IsTrue,  // zero denominator: compile error
    Require<{ NUM <= DEN }>: IsTrue,  // NUM/DEN > 1.0:   compile error
{
    _private: (),
}

// Type aliases for each dimension:
pub type FitnessConst<const N: u64, const D: u64> =
    Metric<{QualityMetricKind::Fitness}, N, D>;
pub type PrecisionConst<const N: u64, const D: u64> =
    Metric<{QualityMetricKind::Precision}, N, D>;
pub type F1Const<const N: u64, const D: u64> =
    Metric<{QualityMetricKind::F1}, N, D>;
pub type GeneralizationConst<const N: u64, const D: u64> =
    Metric<{QualityMetricKind::Generalization}, N, D>;
pub type SimplicityConst<const N: u64, const D: u64> =
    Metric<{QualityMetricKind::Simplicity}, N, D>;
\end{lstlisting}

\subsection{The Metric Lattice}

The five conformance dimensions form a partial order under the interpretation of which
constraints on which. Fitness and precision are primary; F1 aggregates them; generalization
and simplicity are orthogonal secondary dimensions. Figure~\ref{fig:metric-lattice} shows
the metric lattice.

\begin{figure}[ht]
\centering
\begin{tikzpicture}[
  node distance=1.5cm and 2.8cm,
  every node/.style={font=\small},
  mnode/.style={draw, rectangle, rounded corners=2pt, minimum width=2.8cm,
                minimum height=0.85cm, align=center, fill=gray!8},
  >=latex
]
% Nodes
\node[mnode] (fitness)      {\texttt{FitnessConst}\\$\in [0,1]$};
\node[mnode, right=of fitness] (precision) {\texttt{PrecisionConst}\\$\in [0,1]$};
\node[mnode, below=1.2cm of fitness, xshift=1.4cm] (f1) {\texttt{F1Const}\\$= \frac{2 \cdot \text{fit} \cdot \text{prec}}{\text{fit} + \text{prec}}$};
\node[mnode, right=3.5cm of fitness, yshift=-1.6cm] (gen) {\texttt{GeneralizationConst}\\$\in [0,1]$};
\node[mnode, right=3.5cm of fitness, yshift=-3.2cm] (simp) {\texttt{SimplicityConst}\\$\in [0,1]$};
\node[mnode, below=1.8cm of f1, xshift=1.7cm] (verdict) {\texttt{ConformanceVerdict}\\(aggregate)};

% Edges
\draw[->] (fitness)   -- (f1);
\draw[->] (precision) -- (f1);
\draw[->] (f1)        -- (verdict);
\draw[->] (gen)       -- (verdict);
\draw[->] (simp)      -- (verdict);

% Between01 annotation
\node[font=\scriptsize, gray, right=0.3cm of fitness, yshift=0.35cm]
  {$\texttt{Between01}$};

\end{tikzpicture}
\caption{The conformance metric lattice. Fitness and precision contribute to F1;
  generalization and simplicity are independent quality dimensions. All five flow into the
  \texttt{ConformanceVerdict} aggregate. Every metric value is
  \texttt{Between01}-bounded: a value outside $[0, 1]$ is a compile error.}
\label{fig:metric-lattice}
\end{figure}

\begin{lstlisting}[language=Rust,
  caption={Compile-fail receipt: \texttt{FitnessConst<2, 1>} violates
    \texttt{Between01}. A fitness value of $2.0$ is structurally impossible.},
  label={lst:fitness-out-of-range}]
// COMPILE-FAIL: Between01 bound --- fitness value 2/1 > 1.0 is rejected.
// Paper: Carmona et al. (2018) --- fitness in [0, 1].
use wasm4pm_compat::conformance::FitnessConst;

fn main() {
    let _: FitnessConst<2, 1> = FitnessConst::new();
    // ERROR: Require<false>: IsTrue is not satisfied (NUM > DEN)
}
\end{lstlisting}

% -----------------------------------------------------------------------------
\section{New Modules: Process Cube, Temporal, Object Lifecycle}
\label{sec:canon-new}
% -----------------------------------------------------------------------------

\subsection{Process Cube (van der Aalst 2013)}

The process cube framework~\cite{vanderaalst2013processcubes} provides a multi-dimensional
structure for comparing process behavior across different slices of an event log. A process
cube is defined over a set of \emph{dimensions} (axes), and for each combination of
dimension values a separate event-log slice can be projected and mined.

\begin{definition}[Process Cube Dimension]
\label{def:process-cube-dim}
A process cube dimension is a named axis of comparison. In \texttt{src/process\_cube.rs},
a dimension is encoded as \texttt{CubeDimension<NAME>} where \texttt{NAME} is a
\texttt{\&'static str} const parameter. \texttt{CubeDimension<"resource">} and
\texttt{CubeDimension<"time">} are distinct types at compile time.
\end{definition}

\begin{lstlisting}[language=Rust,
  caption={\texttt{CubeDimension<NAME>}: a named process cube axis. The const-string
    parameter makes \texttt{"resource"} and \texttt{"time"} distinct types.},
  label={lst:cube-dimension}]
/// A named dimension in the process cube.
/// CubeDimension<"resource"> and CubeDimension<"time"> are different types.
pub struct CubeDimension<const NAME: &'static str>;

/// A slice: one axis value binding.
pub struct CubeSlice<D, V> {
    /// The dimension this slice cuts along.
    pub dimension: PhantomData<D>,
    /// The concrete value for this dimension.
    pub value: V,
}

/// A cell: a conjunction of dimension-value bindings.
pub struct CubeCell<Slices> {
    pub slices: Slices,
    _private:   PhantomData<()>,
}
\end{lstlisting}

The process cube engine --- slicing event logs per cell, applying discovery per slice,
computing cross-cell conformance distances --- graduates to \texttt{wasm4pm}. The shapes
above travel with the evidence into the engine.

\subsection{Temporal Ordering Law}

The temporal module (\texttt{src/temporal.rs}) encodes the structural vocabulary for
temporal event ordering. Temporal profiles and temporal conformance checking are defined
in Adriansyah et al.~\cite{adriansyah2015measuring}, which extends the conformance
framework with a time dimension. The structural shapes produced by a temporal conformance
check --- temporal ordering relations, sojourn time bounds, temporal deviation records --- are
carried here.

\begin{lstlisting}[language=Rust,
  caption={\texttt{TemporalOrder}: the four canonical temporal relations. Deriving the
    ordering from event timestamps graduates to \texttt{wasm4pm}.},
  label={lst:temporal-order}]
#[derive(Clone, Copy, PartialEq, Eq, Debug, Hash)]
pub enum TemporalOrder {
    Before,      // strictly precedes
    After,       // strictly follows
    Concurrent,  // no strict ordering established
    Unknown,     // ordering not determinable
}

/// Evidence wrapper: a value tagged with a temporal ordering context.
pub struct TemporalEvidence<T> {
    pub value:   T,
    pub context: TemporalContext,
}
\end{lstlisting}

\subsection{Object Lifecycle Const-Generic Phase Transitions}

Object-centric process mining requires reasoning about object lifecycles: the sequence of
phases an object passes through from creation to deletion. In \texttt{src/object\_lifecycle.rs},
object lifecycle phases are encoded as a \texttt{ConstParamTy} enum, and lifecycle
transitions are enforced at the type level.

\begin{definition}[Lawful Lifecycle Transition]
\label{def:lifecycle-transition}
A lifecycle transition $\text{FROM} \to \text{TO}$ is lawful iff it is represented by a
\texttt{LifecycledObject<T, \{FROM\}>} method that returns a
\texttt{LifecycledObject<T, \{TO\}>}. Only the following transitions have method
implementations: $\texttt{Created} \to \texttt{Active}$, $\texttt{Active} \to
\texttt{Modified}$, $\texttt{Active} \to \texttt{Archived}$, $\texttt{Modified} \to
\texttt{Modified}$, $\texttt{Modified} \to \texttt{Archived}$, $\texttt{Archived} \to
\texttt{Deleted}$.
\end{definition}

\begin{lstlisting}[language=Rust,
  caption={Object lifecycle phase transitions enforced at the type level.
    \texttt{LifecycledObject<T, \{Active\}>} has no \texttt{.delete()} method ---
    direct deletion from Active is structurally impossible.},
  label={lst:lifecycle-transitions}]
#[derive(ConstParamTy, Clone, Copy, PartialEq, Eq, Debug, Hash)]
pub enum ObjectLifecyclePhase {
    Created, Active, Modified, Archived, Deleted,
}

pub struct LifecycledObject<T, const PHASE: ObjectLifecyclePhase> {
    pub inner: T,
    _state:    PhantomData<ObjectState<PHASE>>,
}

// Only lawful transitions have methods:
impl<T> LifecycledObject<T, { ObjectLifecyclePhase::Created }> {
    pub fn activate(self)
        -> LifecycledObject<T, { ObjectLifecyclePhase::Active }> { ... }
}
impl<T> LifecycledObject<T, { ObjectLifecyclePhase::Active }> {
    pub fn modify(self)
        -> LifecycledObject<T, { ObjectLifecyclePhase::Modified }> { ... }
    pub fn archive(self)
        -> LifecycledObject<T, { ObjectLifecyclePhase::Archived }> { ... }
}
// No .delete() on Active: that path does not exist.
impl<T> LifecycledObject<T, { ObjectLifecyclePhase::Archived }> {
    pub fn delete(self)
        -> LifecycledObject<T, { ObjectLifecyclePhase::Deleted }> { ... }
}
\end{lstlisting}

\subsection{The E0391 Type Alias Fix}

\noindent\textbf{Technical note: nightly compiler cycle with \texttt{ConstParamTy} enums.}

When a \texttt{ConstParamTy} enum is used as a const generic parameter in a where-bound
expression --- such as \texttt{Require<\{ PHASE == ObjectLifecyclePhase::Active \}>: IsTrue}
--- the nightly Rust compiler (as of the May~2026 manufacturing run) sometimes emits
\texttt{E0391: cycle detected when computing the variance of X}. This happens because the
compiler attempts to evaluate the generic const expression for the where-bound before it
has fully resolved the enum's variance.

The fix is to introduce a type alias that mediates between the enum usage and the
where-bound. Instead of writing the where-bound directly on the parameterized struct,
break the cycle by naming an intermediate alias:

\begin{lstlisting}[language=Rust,
  caption={The E0391 type-alias fix for nightly const-generic cycles with
    \texttt{ConstParamTy} enums. The alias breaks the variance cycle.},
  label={lst:e0391-fix}]
// Problematic --- may trigger E0391 in nightly:
pub struct DirectTransition<T, const FROM: ObjectLifecyclePhase,
                               const TO: ObjectLifecyclePhase>
where
    Require<{ FROM != TO }>: IsTrue,
{ ... }

// Fix: introduce a type alias that avoids the direct where-bound cycle.
// Use the alias at call sites; the alias carries the constraint lazily.
pub type ActiveToModified<T> =
    LifecycledObject<T, { ObjectLifecyclePhase::Modified }>;
pub type ArchivedObject<T> =
    LifecycledObject<T, { ObjectLifecyclePhase::Archived }>;
\end{lstlisting}

This pattern was used across several lifecycle and const-generic modules in the May~2026
manufacturing run to avoid E0391 without disabling the compile-time phase-transition law.
The aliases carry the phase information in their name, preserving readability, while the
compiler's variance inference avoids the cycle by not needing to evaluate the
where-bound expression during type construction.

% -----------------------------------------------------------------------------
\section{The Process Canon Surface Summary}
\label{sec:canon-summary}
% -----------------------------------------------------------------------------

Table~\ref{tab:canon-surface} summarizes the process mining formal objects encoded in this
crate, the paper authority for each, the Rust type that encodes it, and the compile-fail
receipt that proves the structural law is enforced.

\begin{table}[ht]
\centering
\caption{Process mining canon: paper authority, Rust encoding, and type-law receipts.}
\label{tab:canon-surface}
\begin{tabular}{@{}p{2.8cm}p{2.6cm}p{3.2cm}p{3.2cm}@{}}
\toprule
\textbf{Canon object} & \textbf{Paper} & \textbf{Rust type} & \textbf{Law enforced} \\
\midrule
OCEL Event          & OCEL 2.0 \cite{berti2023ocel20}         & \texttt{OcelEvent}                     & Activity + id required \\
OCEL Object         & OCEL 2.0 \cite{berti2023ocel20}         & \texttt{OcelObject}                    & Type + id required \\
E2O Link            & OCEL 2.0 \cite{ghahfarokhi2021ocel}     & \texttt{EventObjectLink}               & Dangling links refused \\
O2O Link            & OCEL 2.0 \cite{ghahfarokhi2021ocel}     & \texttt{ObjectObjectLink}              & Dangling links refused \\
XES Log             & IEEE 1849 \cite{ieee2016xes}             & \texttt{XesLog}                        & OcelLog $\neq$ XesLog \\
Case-centric marker & XES \cite{verbeek2010xes}               & \texttt{CaseCentricMarker}             & No OCEL/XES coercion \\
Petri arc           & Murata 1989 \cite{murata1989petri}      & \texttt{PlaceToTransitionArc}          & P$\to$P forbidden \\
Petri arc           & Murata 1989 \cite{murata1989petri}      & \texttt{TransitionToPlaceArc}          & T$\to$T forbidden \\
WF-net soundness    & van der Aalst 1998 \cite{vanderaalst1998workflow} & \texttt{WfNetConst<SOUNDNESS>} & Non-forgeable witness \\
POWL projection     & Kourani 2024 \cite{kourani2024powl}     & \texttt{ProcessTreeProjectable}        & Irreducible $\neq$ projectable \\
Acyclicity witness  & Kourani 2024 \cite{kourani2024powl}     & \texttt{AcyclicPartialOrder}           & PartialOrder $\not\Rightarrow$ DAG \\
Fitness             & Carmona 2018 \cite{carmona2018conformance} & \texttt{FitnessConst<N,D>}          & $[0,1]$ compile-enforced \\
Precision           & Carmona 2018 \cite{carmona2018conformance} & \texttt{PrecisionConst<N,D>}        & $[0,1]$ compile-enforced \\
F1 score            & Carmona 2018 \cite{carmona2018conformance} & \texttt{F1Const<N,D>}               & $[0,1]$ compile-enforced \\
Generalization      & Carmona 2018 \cite{carmona2018conformance} & \texttt{GeneralizationConst<N,D>}   & $[0,1]$ compile-enforced \\
Simplicity          & Carmona 2018 \cite{carmona2018conformance} & \texttt{SimplicityConst<N,D>}       & $[0,1]$ compile-enforced \\
Process cube dim.   & van der Aalst 2013 \cite{vanderaalst2013processcubes} & \texttt{CubeDimension<NAME>} & Dimension names typed \\
Temporal order      & Adriansyah 2015 \cite{adriansyah2015measuring} & \texttt{TemporalOrder}         & Four-valued relation \\
Object lifecycle    & OCEL 2.0 \cite{berti2023ocel20}         & \texttt{LifecycledObject<T,PHASE>}     & Unlawful transitions absent \\
\bottomrule
\end{tabular}
\end{table}

Every row in this table corresponds to at least one compile-fail fixture in
\texttt{tests/ui/compile\_fail/} and at least one compile-pass fixture in
\texttt{tests/ui/compile\_pass/}. The compile-fail fixture proves that the violation is
rejected; the compile-pass fixture proves that the lawful path remains open. This
bidirectional receipt is the operational definition of a type-law surface in
\texttt{wasm4pm-compat}.

The central invariant of this chapter is:

\begin{quote}
Process mining formal objects are not documentation or convention. They are types. Their
structural laws are not tests. They are compile-time obligations. An application that does
not conform to these obligations does not compile.
\end{quote}
